\documentclass[a4paper,11pt]{article}
\usepackage[utf8]{inputenc}
\usepackage[T1]{fontenc}
\usepackage[spanish]{babel}
\usepackage{geometry}
\usepackage{array}
\usepackage{colortbl}
\usepackage{xcolor}
\usepackage{booktabs}
\usepackage{enumitem}
\usepackage{fancyhdr}
\usepackage{tikz}

\geometry{left=2cm, right=2cm, top=2.5cm, bottom=2cm}

\definecolor{violeta}{HTML}{7B2D8E}
\definecolor{azulg}{HTML}{3498DB}
\definecolor{rojog}{HTML}{E74C3C}
\definecolor{gris}{HTML}{F2F2F2}
\definecolor{verde}{HTML}{27AE60}
\definecolor{naranja}{HTML}{E67E22}

\pagestyle{fancy}
\fancyhf{}
\fancyhead[L]{\small\textbf{Nuevo Compañeros 1 — U03: La Familia}}
\fancyhead[R]{\small\textbf{Píldora formativa 3.5}}
\fancyfoot[C]{\thepage}

\setlength{\parindent}{0pt}

\begin{document}

\begin{center}
{\LARGE\textbf{PÍLDORA FORMATIVA 3.5}}\\[0.3cm]
{\Large Comprensión comunicativa: texto mapeado del diálogo}\\[0.2cm]
{\normalsize Sección Comunicación (p.38) — Proyectable — 6 diapositivas}\\[0.2cm]
{\small Versión enriquecida secuencial para Fases C1--C4}
\end{center}

\vspace{0.5cm}

\noindent\fbox{\parbox{\dimexpr\linewidth-2\fboxsep-2\fboxrule}{%
\textbf{Contenido para el profesor}\\[0.2cm]
El diálogo En el recreo (p.38) es el texto comunicativo central de la unidad. Cuatro personajes (Pablo, Graciela, Julia, Jorge) hablan de sus familias durante el recreo. El diálogo integra vocabulario de parentesco ya adquirido, verbo \textit{tener} en contexto funcional, interrogativos ya formalizados y funciones comunicativas clave: preguntar por hermanos, responder sobre familia, invitar, aceptar, rechazar.\\[0.2cm]
El texto tiene 7 huecos que los estudiantes completan con palabras del recuadro: formas verbales (\textit{tienes, tengo, Tienen}) e interrogativos (\textit{Cuántos, Cómo, Dónde}) más el verbo \textit{estar} (\textit{Está}).
}}

\vspace{0.8cm}

%% ============================================================
%% DIAPOSITIVA 1
%% ============================================================

\noindent\textcolor{violeta}{\rule{\linewidth}{2.5pt}}

\section*{DIAPOSITIVA 1 — ¿QUÉ PASA EN EL RECREO?}

{\small Fase C1 (Contextualización) + Pre-visionado — Técnica: Activación de esquema previo}\\
{\small\textit{Principio:} Antes de ver el vídeo, el estudiante activa su conocimiento previo sobre la situación comunicativa (el recreo). La pregunta guía reduce la carga cognitiva al dar un foco atencional único.}

\vspace{0.5cm}

\subsection*{En pantalla}

\begin{center}
{\Large\textbf{EN EL RECREO}}
\end{center}

\vspace{0.3cm}

Fotograma congelado del vídeo: los cuatro personajes en el patio.\\
Debajo: \textit{¿Qué veis? ¿Dónde están? ¿De qué van a hablar?}

\vspace{0.3cm}

Cuatro fotos de los personajes con sus nombres: \textbf{Pablo}, \textbf{Graciela}, \textbf{Julia}, \textbf{Jorge}.

\vspace{0.3cm}

\noindent Nube de palabras clave:

\begin{center}
\begin{tikzpicture}
\node[font=\large\bfseries, text=violeta] at (0,0) {familia};
\node[font=\normalsize, text=azulg] at (2.5,0.3) {hermano};
\node[font=\normalsize, text=rojog] at (5,0) {hermana};
\node[font=\small, text=verde] at (1,-0.6) {primo};
\node[font=\large, text=naranja] at (3.5,-0.5) {gemelas};
\node[font=\small, text=violeta] at (5.5,-0.7) {hija única};
\node[font=\normalsize, text=azulg] at (0.5,-1.2) {recreo};
\node[font=\small, text=rojog] at (3,-1.3) {merendar};
\end{tikzpicture}
\end{center}

\vspace{0.3cm}

Instrucción: \textit{Van a hablar de sus familias. La primera vez, solo escuchad y responded:}

\begin{center}
{\large\textbf{¿Quién tiene hermanos? ¿Quién no tiene hermanos?}}
\end{center}

\vspace{0.3cm}

\begin{center}
\renewcommand{\arraystretch}{1.5}
\begin{tabular}{|>{\centering\arraybackslash}p{6cm}|>{\centering\arraybackslash}p{6cm}|}
\hline
\rowcolor{verde!15}
\textbf{TIENE HERMANOS} & \textbf{NO TIENE HERMANOS} \\
\hline
 & \\[1cm]
\hline
\end{tabular}
\end{center}

\vspace{0.5cm}

\subsection*{Instrucciones para el profesor}

\begin{enumerate}[leftmargin=1.5cm, itemsep=4pt]
\item Proyecte el fotograma congelado. Pregunte: \textit{¿Qué veis? ¿Dónde están estos chicos?} Recoja descripciones breves. Pregunte: \textit{¿De qué creéis que van a hablar?} Anote 2--3 hipótesis en la pizarra.
\item Señale las fotos de los personajes. Diga sus nombres. Pida que los repitan.
\item Señale la nube de palabras. Lea las palabras en voz alta. Para las palabras nuevas (\textit{gemelas, hija única, merendar}), use gestos o dibujos rápidos. Pregunte: \textit{¿Cuáles conocéis ya? ¿Cuáles son nuevas?}
\item Plantee la pregunta guía en voz alta.
\item Reproduzca el vídeo (o audio pista 36) SIN pausas.
\item Al terminar, pregunte: \textit{¿Quién tiene hermanos?} Recoja respuestas a mano alzada. Vuelva al fotograma: \textit{¿Se cumplen vuestras hipótesis?}
\item Respuesta esperada: Pablo (dos hermanas gemelas), Graciela (una hermana). NO tienen: Julia (hija única), Jorge (no hermanos, pero primos).
\end{enumerate}

\newpage

%% ============================================================
%% DIAPOSITIVA 2
%% ============================================================

\noindent\textcolor{violeta}{\rule{\linewidth}{2.5pt}}

\section*{DIAPOSITIVA 2 — COMPLETAD EL DIÁLOGO}

{\small Fase C2 (Comprensión del modelo) + Durante el visionado — Técnica: Segundo visionado con tarea detallada}\\
{\small\textit{Principio:} El segundo visionado con los huecos del libro obliga al estudiante a procesar formas específicas (verbos e interrogativos). El recuadro de opciones reduce la carga cognitiva (reconocimiento vs. recuerdo libre).}

\vspace{0.5cm}

\subsection*{En pantalla}

Mini-tabla de observación:

\vspace{0.3cm}

\begin{center}
\renewcommand{\arraystretch}{1.5}
\begin{tabular}{|>{\centering\arraybackslash}p{3cm}|>{\centering\arraybackslash}p{5cm}|>{\centering\arraybackslash}p{5cm}|}
\hline
\rowcolor{azulg!15}
\textbf{Personaje} & \textbf{¿Tiene hermanos? ¿Cuántos?} & \textbf{Palabra del recuadro usada} \\
\hline
Pablo & & \\
\hline
Graciela & & \\
\hline
Julia & & \\
\hline
Jorge & & \\
\hline
\end{tabular}
\end{center}

\vspace{0.3cm}

El diálogo con los 7 huecos numerados (tal como aparece en el libro, p.38).

\vspace{0.3cm}

\noindent\colorbox{gris}{\parbox{\dimexpr\linewidth-2\fboxsep}{%
\centering\textbf{Recuadro de palabras:} tienes \quad tengo \quad Cuántos \quad Tienen \quad Cómo \quad Dónde \quad Está
}}

\vspace{0.3cm}

Instrucción: \textit{Escuchad otra vez. Completad los huecos con las palabras del recuadro.}

\vspace{0.5cm}

\subsection*{Instrucciones para el profesor}

\begin{enumerate}[leftmargin=1.5cm, itemsep=4pt]
\item Pida que abran el libro en p.38 y lean las palabras del recuadro.
\item Pregunte: \textit{¿Qué tipo de palabras son? ¿Verbos? ¿Preguntas?} Los estudiantes deben notar que hay formas de \textit{tener} (tienes, tengo, Tienen), interrogativos (Cuántos, Cómo, Dónde) y una forma nueva (Está).
\item Señale la mini-tabla de observación. Explique: \textit{Mientras veis el vídeo, anotad también cuántos hermanos tiene cada personaje.}
\item Reproduzca el vídeo (o audio) por segunda vez. Los estudiantes completan los huecos y la tabla mientras escuchan.
\item Pause si es necesario, especialmente antes de cada hueco.
\item Tras la escucha, pida que comparen en parejas (1 minuto): primero los huecos, después la tabla.
\end{enumerate}

\vspace{0.3cm}

\noindent\textbf{Respuestas:} 1=tienes, 2=tengo, 3=Cuántos, 4=Tienen, 5=Cómo, 6=Dónde, 7=Está.

\newpage

%% ============================================================
%% DIAPOSITIVA 3
%% ============================================================

\noindent\textcolor{violeta}{\rule{\linewidth}{2.5pt}}

\section*{DIAPOSITIVA 3 — ¿QUIÉN ES QUIÉN?}

{\small Fase C2 (Comprensión profunda) + Post-visionado — Técnica: Comprensión escalonada global → detalle}\\
{\small\textit{Principio:} Tras completar el diálogo, el estudiante demuestra comprensión profunda completando frases con los nombres correctos (Act. 2) y evaluando afirmaciones V/F (Act. 3).}

\vspace{0.5cm}

\subsection*{En pantalla}

\noindent\textbf{Sección A — Completad con los nombres (Act. 2):}
\begin{enumerate}[leftmargin=1.5cm, itemsep=3pt]
\item \rule{3cm}{0.4pt} tiene primos en Argentina.
\item \rule{3cm}{0.4pt} es hija única.
\item La hermana de \rule{3cm}{0.4pt} se llama Rebeca.
\item \rule{3cm}{0.4pt} es dos años mayor que Graciela.
\item Las hermanas de \rule{3cm}{0.4pt} son gemelas.
\item \rule{3cm}{0.4pt} y \rule{3cm}{0.4pt} tienen ocho años.
\end{enumerate}

\vspace{0.3cm}

\noindent\textbf{Sección B — ¿Verdadero o falso? (Act. 3):}
\begin{enumerate}[leftmargin=1.5cm, itemsep=3pt]
\item Los gemelos son dos hermanos que nacen el mismo día. (\quad)
\item Los gemelos son dos hermanos iguales. (\quad)
\item Los gemelos son siempre del mismo sexo. (\quad)
\end{enumerate}

\vspace{0.3cm}

\noindent\textbf{Sección C — Mapa de relaciones:}

\begin{center}
\renewcommand{\arraystretch}{1.5}
\begin{tabular}{|>{\centering\arraybackslash}p{3cm}|>{\centering\arraybackslash}p{3cm}|>{\centering\arraybackslash}p{5cm}|}
\hline
\rowcolor{violeta!15}
\textbf{Personaje A} & \textbf{Personaje B} & \textbf{Relación} \\
\hline
Pablo & Graciela & compañeros \\
\hline
Pablo & Julia & \rule{3cm}{0.4pt} \\
\hline
Pablo & Jorge & \rule{3cm}{0.4pt} \\
\hline
Graciela & Jorge & \rule{3cm}{0.4pt} \\
\hline
\end{tabular}
\end{center}

Instrucción: \textit{Dibujad flechas entre los personajes. Escribid qué relación tienen.}

\vspace{0.5cm}

\subsection*{Instrucciones para el profesor}

\begin{enumerate}[leftmargin=1.5cm, itemsep=4pt]
\item Pida que completen la Sección A individualmente con el diálogo abierto (2 min).
\item Corrección en parejas. Puesta en común selectiva: solo los ítems donde hubo desacuerdo.
\item Respuestas Act. 2: 1. Jorge, 2. Julia, 3. Graciela, 4. Rebeca, 5. Pablo, 6. Adela y Luisa.
\item Para la Sección B, lea cada afirmación en voz alta. Pida pulgar arriba (V) o pulgar abajo (F).
\item Respuestas Act. 3: 1. V, 2. V, 3. F.
\item Para el ítem 3 (F), pregunte: \textit{¿Conocéis gemelos de diferente sexo?} Aclare que en español, \textit{gemelos} se usa para ambos casos.
\item Para la Sección C, pida que dibujen las conexiones en su cuaderno (2 min). Muestre un ejemplo en la pizarra: Pablo $\leftrightarrow$ Adela y Luisa = hermanas gemelas.
\end{enumerate}

\newpage

%% ============================================================
%% DIAPOSITIVA 4
%% ============================================================

\noindent\textcolor{violeta}{\rule{\linewidth}{2.5pt}}

\section*{DIAPOSITIVA 4 — ¿CÓMO PREGUNTAN? ¿CÓMO RESPONDEN?}

{\small Fase C2--C3 (Extracción de funciones + Post-visionado) — Técnica: Análisis de funciones comunicativas con mapa en pizarra}\\
{\small\textit{Principio:} El estudiante comprende el contenido del diálogo y ahora se extraen las funciones comunicativas explícitas. El profesor guía la extracción y escribe la función en la pizarra. Se construye el esquema comunicativo de forma inductiva.}

\vspace{0.5cm}

\subsection*{En pantalla}

Tabla con dos columnas (las filas se revelan una a una con clic):

\vspace{0.3cm}

\begin{center}
\renewcommand{\arraystretch}{1.4}
\begin{tabular}{|>{\raggedright\arraybackslash}p{4.5cm}|>{\raggedright\arraybackslash}p{8.5cm}|}
\hline
\rowcolor{violeta!15}
\textbf{FUNCIÓN} & \textbf{EJEMPLO DEL DIÁLOGO} \\
\hline
Preguntar si tiene hermanos & \textit{¿Tienes hermanos?} \\
\hline
Decir que no tiene hermanos & \textit{No, soy hija única.} \\
\hline
Describir su familia & \textit{Tengo una hermana dos años mayor que yo.} \\
\hline
Preguntar por la edad & \textit{¿Cuántos años tienen?} \\
\hline
Decir la edad & \textit{Tienen 8 años.} \\
\hline
Preguntar por el nombre & \textit{¿Cómo se llaman?} \\
\hline
Decir el nombre & \textit{Se llaman Adela y Luisa.} \\
\hline
Invitar & \textit{Podéis venir a mi casa a merendar.} \\
\hline
Aceptar & \textit{¡Genial!} \\
\hline
Rechazar con disculpa & \textit{Yo no puedo, lo siento.} \\
\hline
Preguntar dónde vive & \textit{¿Dónde vives?} \\
\hline
Decir dónde vive & \textit{Vivo en la calle del instituto.} \\
\hline
Indicar proximidad & \textit{Está muy cerca de aquí.} \\
\hline
\end{tabular}
\end{center}

\vspace{0.5cm}

\subsection*{Instrucciones para el profesor}

\begin{enumerate}[leftmargin=1.5cm, itemsep=4pt]
\item Texto mapeado con marcador de BlinkLearning. Escriba en la pizarra tres funciones: preguntar, responder/describir, invitar/aceptar/rechazar. Marque cada función con un color diferente.
\item Pregunte: \textit{¿Qué puede significar cada color?} Pida que relacionen el color con la función.
\item Para cada función, pregunte: \textit{¿Cómo pregunta Graciela por los hermanos?} Los estudiantes buscan en el diálogo y leen la frase. Haga clic para revelar la fila correspondiente.
\item Al llegar a \textit{Está muy cerca de aquí}, señale que es una forma nueva: \textit{Esta es del verbo estar. Significa 'is located'. Lo usamos para decir dónde está algo. Ahora solo necesitamos esta forma.}
\item Proponga frases nuevas para clasificar por función e invítelos a inventar las suyas.
\item Tras revelar todas las funciones, reparta la tarjeta Esquema comunicativo: hablar de la familia (Caja 2). Pida que comparen con lo que acaban de descubrir.
\item Ampliación opcional (T32 Entrevista al personaje): un estudiante asume el rol de Pablo. El resto le hace preguntas usando las funciones del cuadro.
\end{enumerate}

\newpage

%% ============================================================
%% DIAPOSITIVA 5
%% ============================================================

\noindent\textcolor{violeta}{\rule{\linewidth}{2.5pt}}

\section*{DIAPOSITIVA 5 — PRACTICAD EL DIÁLOGO}

{\small Fase C3 (Práctica guiada) — Técnica: Reproducción del diálogo en parejas con weaning off}\\
{\small\textit{Principio:} Primera producción oral. Los estudiantes reproducen el diálogo en parejas, primero leyendo y después de memoria. El weaning off es explícito: ronda 1 con libro, ronda 2 con esquema, ronda 3 sin nada.}

\vspace{0.5cm}

\subsection*{En pantalla}

Instrucciones paso a paso:

\vspace{0.3cm}

\begin{center}
\renewcommand{\arraystretch}{1.8}
\begin{tabular}{|>{\centering\arraybackslash}p{2.5cm}|>{\raggedright\arraybackslash}p{10cm}|}
\hline
\rowcolor{verde!15}
\textbf{Paso} & \textbf{Instrucción} \\
\hline
\textbf{Shadowing} & Reproducid el vídeo una vez más. Repetid EN VOZ BAJA al mismo tiempo que los personajes, imitando su entonación. \\
\hline
\textbf{Ronda 1} & Con libro abierto. En parejas. Uno es Pablo/Julia, otro es Graciela/Jorge. Leed el diálogo completo. \\
\hline
\textbf{Ronda 2} & Con esquema. Cerrad el libro. Usad solo la tarjeta Esquema comunicativo. Intentad reproducir las funciones con vuestras propias palabras. \\
\hline
\textbf{Ronda 3} & Sin apoyo. Cerrad todo. Intentad recordar las funciones principales. \\
\hline
\end{tabular}
\end{center}

\vspace{0.5cm}

\subsection*{Instrucciones para el profesor}

\begin{enumerate}[leftmargin=1.5cm, itemsep=4pt]
\item Antes de la práctica en parejas, reproduzca el vídeo una vez más para shadowing: los estudiantes repiten en voz baja al mismo tiempo que los personajes ($\sim$1 min).
\item Forme parejas. Asigne roles: A = Pablo/Julia, B = Graciela/Jorge.
\item Ronda 1 (2 min): leen el diálogo en voz alta con el libro abierto. Circule y escuche la pronunciación.
\item Ronda 2 (2 min): cierran el libro, dejan visible la tarjeta Esquema comunicativo. No tiene que ser idéntico al diálogo original.
\item Ronda 3 (1 min): sin ningún apoyo. Solo las funciones que recuerden.
\item Pida a 1--2 parejas voluntarias que representen ante la clase.
\end{enumerate}

\newpage

%% ============================================================
%% DIAPOSITIVA 6
%% ============================================================

\noindent\textcolor{violeta}{\rule{\linewidth}{2.5pt}}

\section*{DIAPOSITIVA 6 — CREAD VUESTRO DIÁLOGO}

{\small Fase C4 (Producción autónoma) — Técnica: Producción semi-libre con tarjetas de estrategia}\\
{\small\textit{Principio:} El estudiante crea su propio diálogo con datos personales reales. Las tarjetas de estrategia (Comodín, Tiempo Extra) reducen la ansiedad en la primera producción semi-libre de la unidad.}

\vspace{0.5cm}

\subsection*{En pantalla}

\begin{center}
{\Large\textbf{CREA TU PROPIO DIÁLOGO}}
\end{center}

\vspace{0.3cm}

Pasos (de la Act. 4 del libro):
\begin{enumerate}[leftmargin=1.5cm, itemsep=3pt]
\item Busca fotos de tu familia (o dibuja / describe de memoria).
\item Prepara un diálogo con un compañero. Usad las funciones:
  \begin{itemize}[leftmargin=1cm]
  \item Presentar a alguien
  \item Hablar de la familia
  \item Preguntar y decir la edad, el nombre, dónde vivís
  \end{itemize}
\item Practicad el diálogo varias veces.
\item Representad el diálogo delante de la clase.
\end{enumerate}

\vspace{0.3cm}

Tarjetas de estrategia:

\begin{center}
\begin{tikzpicture}
\node[draw, rounded corners, fill=naranja!20, minimum width=5cm, minimum height=1.5cm, text width=4.5cm, align=center] at (0,0) {\textbf{COMODÍN}\\Puedes pedir UNA palabra al profesor.};
\node[draw, rounded corners, fill=azulg!20, minimum width=5cm, minimum height=1.5cm, text width=4.5cm, align=center] at (7,0) {\textbf{TIEMPO EXTRA}\\Puedes pedir 30 segundos de preparación antes de hablar.};
\end{tikzpicture}
\end{center}

\vspace{0.5cm}

\subsection*{Instrucciones para el profesor}

\begin{enumerate}[leftmargin=1.5cm, itemsep=4pt]
\item Lea los pasos en voz alta.
\item Señale que pueden usar fotos reales de su teléfono o simplemente hablar de su familia de memoria.
\item Tiempo de preparación escrita: 3 min. Circule y ayude con vocabulario.
\item Tiempo de práctica oral en parejas: 2 min.
\item Representación ante la clase: pida voluntarios. Si nadie se ofrece, pregunte: \textit{¿Quién quiere usar el Comodín?} Esto suele animar a participar.
\item Feedback: use recasts (remodele la frase correcta sin señalar el error). Por ejemplo, si un estudiante dice \textit{Mi hermana tiene quince ano}, usted repite: \textit{Ah, tu hermana tiene quince años. ¡Genial!}
\item Esta actividad aplica la técnica T70 (Diálogo paralelo): misma estructura funcional del vídeo original pero con datos personales reales.
\end{enumerate}

\end{document}
